\scanpage{64}
\begin{flushright}atkin page 8\end{flushright}\mspage{8}

Finally we may extend Conjecture 2 in the same way as Conjecture \handwritten{3} extends Conjecture 1. Our proofs(where available) have two difficulties; the first,and most crucial,is to obtain the analogy of (H4q) and show that congruentially we can\inserted{for all $\handwritten{\underline}{\alpha}$} obtain a finite-dimensional subspace of $\handwritten{\underline}{C}$,which we can only do when $\handwritten{\underline}{\handwritten{\Gamma}}_{\handwritten{0}}(\handwritten{\underline}{q})$ has genus zero by the methods of Atkin and O'Brien\handwritten{\uwave}{(3)},where the case $\handwritten{\underline}{q}=13$ of Conjecture 1 is stated. For the remaining values of $\handwritten{\underline}{q}$ we are limited to a few fixed values of $\handwritten{\underline}{\alpha}$. The second difficulty,which only arises for $\handwritten{\underline}{q}\handwritten{\geqslant}29$ when the ``congruence space'' has dimension greater than one,is to obtain a criterion for simultaneous diagonalization such as (H3q). In fact we can use an (H3q*) analogous to (H3*) in \handwritten{\S}2 to the effect that at least one matrix has distinct eigenvalues modulo $\handwritten{\underline}{q}$(it will a fortiori then have distinct eigenvalues modulo $\handwritten{\underline}{q}^{\handwritten{\underline}{\alpha}}$). As indicated in \handwritten{\S}2,there is no problem with (H1q) and (H2q) which carry over for $\handwritten{\underline}{C}$ exactly as for $\handwritten{\underline}{C}^{\handwritten{0}}$.

\handwritten{\uline}{References}.

1. A.O.L.Atkin,Proof of a conjecture of Ramanujan,Glasgow Math.J. \handwritten{\uwave}{8}(1967),14-32.

2. A.O.L.Atkin,Congruences for modular forms,Proceedings of IBM Conference on computers in mathematical research,Blaricum 1966,published North Holland(1967).

3.A.O.L.Atkin and J.N.O'Brien,Some properties of $p(n)$ and $c(n)$ modulo powers of 13,Trans. Amer. Math. Soc. \handwritten{\uwave}{126}(1967),442-459.

4. R.C.Gunning,Lectures on modular forms,Princeton 1962.

5. J.Lehner,Divisibility properties of the Fourier coefficients of the modular invariant $j(\handwritten{\underline}{\handwritten{\tau}})$,Amer. J.Math.\handwritten{\uwave}{71}(1949),\deleted{\unsure{1}\illegible{5}\unsure{5|3}}\inserted{\handwritten{136-148}}.\reversemarginpar\margin{\handwritten{\unsure{tau|Fin|Fiu}}}\normalmarginpar

6. J.Lehner,Further congruence properties of the Fourier coefficients of the modular invariant $j(\handwritten{\underline}{\handwritten{\tau}})$,Amer.J.Math.\handwritten{\uwave}{71}(1949),373-386.

7.J.Lehner,A short course in automorphic functions,New York 1965.

8. R.A.Rankin,Elliptic \deleted{\unsure{M|m}}\handwritten{m}odular \deleted{\unsure{F|f}}\handwritten{f}unctions and \deleted{\unsure{F|f}}\handwritten{f}orms,Lecture notes ,Bloomington,Indiana,1964.

\deleted{\unsure{8|0}}\unsure{9|6|8}. R.A.Rankin and J.M.Rushforth,The coefficients of certain integral modular forms,Proc. Cambridge Phil.Soc. \handwritten{\uwave}{50}(1954),305-308.\margin{\handwritten{\#}}

\unsure{S.R.C.|The} \unsure{Atlas} \unsure{Computer} Laboratory,Chilton,Didcot,England:
% archive sheet number: 63
